\section{Selmer matrices and three-state relations}
\label{sec:lm:local-matrices}

Assume that $V=E[2]$ is irreducible, and write
$F=\Q(V)$ and $G=\Gal(F/\Q)\in\{C_3,S_3\}$.
The symbol rules will be useful because a unit symbol bounds the
$2$-Selmer dimension. We now express local Selmer conditions by a
binary matrix and make that bound constrain the number of singular
blocks. The second half of the section develops the elementary matrix
operations needed to apply this constraint.

Write the Weil pairing additively, and choose a basis of $V$ in which
its matrix is
\[
 J=\begin{pmatrix}0&1\\1&0\end{pmatrix}.
\]
Fix a finite set $S$ containing $2$, $3$, the bad primes, all auxiliary
exclusions, and rational primes whose prime ideals generate the class
groups of every subfield of $F$. Local specifications at $S$ will also
include the archimedean places. The detecting prime $\ell$ for an odd
coefficient system is kept separate: all generators below have valuation
zero there. This separation keeps the set of support types finite as
$\ell$ varies. We shall apply the construction to a fixed curve, after
absorbing any initial quadratic twist into $E$.

\subsection{Global classes and local equations}

At a split prime $p$, choose a place of $F$ and put $H_p=1$.
At a simple prime, choose a place with Frobenius a specified
transposition $c$ and put $H_p=\langle c\rangle$.
Set
\[
 F'_p=F^{H_p},\qquad V_p=V^{H_p}.
\]
Thus $V_p$ has dimension two or one, respectively. The chosen place
lies above a degree-one prime $P_p$ of $F'_p$. By the choice of $S$,
there is $\pi_p\in(F'_p)^\times$ whose divisor away from $S$ is $P_p$.
For $x\in V_p$ define
\begin{equation}\label{eq:lm:generator-class}
 C_p(x)=\cor_{F'_p/\Q}(\chi_{\pi_p}x)\in H^1(\Q,V).
\end{equation}
Here $\chi_{\pi_p}$ is the quadratic Kummer character; its product with
$x$ is a cocycle because $H_p$ fixes $x$.

The inertia coordinate of $C_p(x)$ at $p$ is $x$, and all its other
inertia coordinates outside $S$ vanish. Indeed, local
restriction--corestriction has exactly one ramified summand, at the
degree-one place where $\pi_p$ has valuation one. These coordinates
therefore make the classes for distinct $p$ independent. A root prime
has $V^{\Fr_p}=0$ and contributes no coordinates.

Let $h$ be a twist containing the active nonroot primes in its
discriminant. At a good odd prime $p$ ramified in $h$, inertia identifies
the local Kummer condition with $V_p$. To see surjectivity, halve a
Frobenius-fixed two-torsion point over the maximal unramified extension.
The ramified twisting involution negates its half, so the resulting
Kummer inertia is that two-torsion point. The Kummer condition has
dimension $\dim V_p$ by local duality, giving the asserted isomorphism.

Fix arithmetic Frobenius lifts $\phi_p$. A second local coordinate is
unramified evaluation at $\phi_p$, modulo Frobenius coboundaries.
Let $K_p(x)(\phi_p)$ be this coordinate for the Kummer lift of inertia
$x$. It is determined by $E[4]$ and the local twisting squareclass.
Define the block matrix $B$ by
\begin{equation}\label{eq:lm:raw-matrix}
 B_{qp}(y,x)=\langle C_p(x)(\phi_q),y\rangle
 -\mathbf1_{p=q}\langle K_p(x)(\phi_p),y\rangle,
 \qquad x\in V_p,\quad y\in V_q.
\end{equation}
Thus $Bx=0$ is exactly the collection of active-prime Selmer equations
for $\sum_p C_p(x_p)$. Changing the unramified unit part of the twist
adds $J$ to a split diagonal block. On a simple line it changes nothing,
since the Weil pairing restricts to zero on a one-dimensional space.

Record in the \emph{type} $\mathfrak h_p$ the subgroup $H_p$ and the
squareclasses of $\pi_p$ at the places over $S$, including signatures.
Only finitely many types occur. Finer generator neighborhoods will be
used to reproduce self entries, but they are not part of this type
alphabet and do not enter its codimension cost. Global reciprocity gives
\begin{equation}\label{eq:lm:symmetry-defect}
 B_{qp}+B_{pq}^{t}=b(\mathfrak h_q,\mathfrak h_p),
\end{equation}
where $b$ is the sum of the local cup pairings at $S$.
Indeed, at an active prime decompose each class into its Kummer lift
and its unramified difference. Both planes are isotropic, and the
cross pairing is unramified evaluation paired with inertia. The terms
at $p$ and $q$ give the left side; the sum over all places is zero.
Elsewhere outside $S$ the classes are unramified, so the local pairing
vanishes. This includes $\ell$, where the generators have unit valuation.
The same calculation shows that $b(\mathfrak h,\mathfrak h)$ is alternating.

\subsection{Prescribing the matrix entries}

We next prove that the off-diagonal entries and the rational Legendre
bits can be varied independently. The latter bits will control split
diagonals without disturbing the other blocks. The class-field inputs
are ray-class reciprocity, the Hilbert product formula, and Chebotarev
\citep[Chapter~VI, Sections~7--8, and Chapter~VII, Section~13]{Neukirch1999};
see also \citet[Chapter~V, Theorems~3.23, 5.3 and~5.5,
and Example~5.4]{MilneCFT2020}.

\begin{lemma}[Realization of finite matrix prescriptions]
\label{lem:lm:arithmetic-realization}
Fix finitely many individually admissible oriented prime templates,
including generator neighborhoods at $S$, any additional finite places,
and oriented Frobenius on $E[4]$ and fixed quadratic-radical extensions
unramified outside that support. At arbitrarily fresh primes these
templates admit the following simultaneous prescriptions.
\begin{enumerate}[label=\textup{(\roman*)}]
\item Between distinct vertices, one direction of the double-coset
residue bits is free and reciprocity determines the reverse direction.
The available matrix blocks are all binary matrices for two split
vertices in the $S_3$ case, and $J\F_4$ in the $C_3$ case, with the
rotation algebra identified with $\F_4$. For split--simple vertices
the full vector space is available, and for simple--simple vertices
the scalar entry is available.
\item The rational Legendre bit between the absolute primes is
independent of the matrix block, subject to rational reciprocity in
the reverse direction.
\item Repeating a template fixes its simple self entry and its split
self block modulo a scalar multiple of $J$. A sufficiently primary
split template with trivial $E[4]$-Frobenius has self block in $\F_2J$.
\end{enumerate}
Here \emph{primary} means sufficiently close to $1$ at the finite
places of $S$ and totally positive at the real places. New cross-bit
prescriptions are made outside the support. All assertions concern
finite configurations; an additional moving prime $\ell$ may be
included among the prescribed places.
\end{lemma}

\begin{proof}
\emph{Off-diagonal entries.}
Restriction--corestriction in Equation~\eqref{eq:lm:generator-class}
gives, for distinct vertices,
\begin{equation}\label{eq:lm:double-coset}
 B_{qp}(y,x)=
 \sum_{g\in H_q\backslash G/H_p}
 e_{pq}(g)\langle gx,y\rangle.
\end{equation}
The bit $e_{pq}(g)$ is the quadratic residue of $g\pi_p$ at the
oriented $q$-place in $F^{H_q\cap gH_pg^{-1}}$. Indeed, pair local
transfer with the $H_q$-invariant vector $y$; the orbit length
$[H_q:H_q\cap gH_pg^{-1}]$ is precisely the residue degree in the
norm calculation. In this intermediate field, each of $\pi_q$ and
$g\pi_p$ has its off-support divisor at a single prime: its oriented
Frobenius generates the subgroup fixing the generator. Hilbert
reciprocity therefore relates $e_{pq}(g)$ to $e_{qp}(g^{-1})$, with
correction only at $S$.

For a new vertex the independent forward bits are residue conditions
at distinct primes of its own field $F'_p$. Move an initial template
away from the old primes and choose a degree-one prime in the ray class
whose quotient generator has the prescribed neighborhoods at $S$,
real signs, and nonzero residues at those old primes. Multiplying the
template generator by this quotient gives the required new generator.
These conditions also preserve the admissible fixed Frobenius data:
changes by local units at old primes are trivial on the intersection
with extensions unramified outside the fixed support. Chebotarev thus
imposes the ray and fixed-label conditions simultaneously. Prime ideals
of degree greater than one over $\Q$ have density zero among prime
ideals of $F'_p$, so requiring degree one loses no admissible class.
The same argument permits additional finite support, including $\ell$.
Reverse bits are then determined by reciprocity. Rational norm bits at
already chosen root primes may be imposed as additional residue
conditions in the ray modulus: finite-field norm maps are surjective,
and these primes contribute no matrix coordinates.

\emph{Matrix span and the norm bit.}
The norm to $\Q$ writes the rational absolute-prime Legendre bit as
the sum of the double-coset bits, plus a fixed support-unit correction.
For two split vertices the matrices in
Equation~\eqref{eq:lm:double-coset} are $Jg$. They span
$\operatorname{Mat}_2(\F_2)$ for $S_3$ and $J\F_4$ for $C_3$.
A nontrivial rotation satisfies $1+R+R^2=0$: changing these three
bits fixes the matrix and changes their sum. For split--simple
vertices, the three nonzero vectors or dual vectors likewise span
the full space and have the same odd relation. For two simple vertices,
the matching double coset has zero pairing and the other has nonzero
pairing. The first bit changes the norm bit alone. This proves
\textup{(i)} and \textup{(ii)}.

\emph{Self entries.}
The nonidentity double cosets in self-transfer are unramified at the
oriented prime. The identity term is a scalar times the pairing and
is therefore zero on a simple line. At a split vertex the two rotation
terms have fixed sum by reciprocity; changing both alters the block
only by a scalar multiple of $J$. Every other term has an involution
representative $\tau$ exchanging the oriented place with another one.
For a simple vertex this includes a transposition other than $c$
representing the nonmatching self double coset.

Work in $F^\tau$ and put
\[
 \alpha=\pi_p\tau\pi_p,\qquad
 \beta=\pi_p+\tau\pi_p,\qquad
 F=F^\tau(\sqrt{D_\tau}).
\]
The distinct divisors of $\pi_p$ and $\tau\pi_p$ ensure $\beta\ne0$.
If $v_0$ is the place below the oriented prime, the residue symbol of
$\tau\pi_p$ there is $(\alpha,\beta)_{v_0}$: $\alpha$ has valuation
one and $\beta$ is a unit with the required residue. Outside
$S\cup\{v_0\}$, $\alpha$ is a unit. When $\beta$ is a unit, both
$(\alpha,\beta)_v$ and $(-D_\tau,\beta)_v$ vanish. When $\beta$ has
positive valuation, put
$\gamma=(\pi_p-\tau\pi_p)/\sqrt{D_\tau}\in F^\tau$. The identity
\[
 \beta^2-4\alpha=D_\tau\gamma^2
\]
shows that $\alpha$ has residue squareclass $-D_\tau$; the two
symbols are again equal. At $v_0$, $(-D_\tau,\beta)_{v_0}=0$.
Reciprocity applied to both symbols gives, in additive notation,
\begin{equation}\label{eq:lm:self-spin}
 (\alpha,\beta)_{v_0}
 =\sum_{v\in S}\bigl((\alpha,\beta)_v+(-D_\tau,\beta)_v\bigr).
\end{equation}

At finite support places the right side is locally constant in a
sufficiently small neighborhood of the template. At a real place
where $F/F^\tau$ is complex, $\alpha$ and $-D_\tau$ are positive.
At a split real place, opposite signs of $\pi_p$ and $\tau\pi_p$
make the two terms cancel; equal signs determine their sum. Thus
signatures suffice, without archimedean size restrictions.
For a primary template, $\alpha$ is a square at the finite support
places and $\beta$ has the squareclass of $2$. The support sum becomes
$\sum_{v\in S}(-D_\tau,2)_v=0$, since the terms outside $S$ vanish.
The involution terms are therefore zero. Together with the local
Kummer lift determined by $E[4]$, this proves \textup{(iii)}.
\end{proof}

\subsection{Symmetrization and the Selmer bound}

The realization lemma supplies the allowable matrix blocks. We now remove
their fixed-support reciprocity defect and compare the resulting kernels
with actual Selmer classes. The fixed support types, stationary
environment and additional local conditions will account for the loss
in dimension.

Choose a matrix rule $L(\mathfrak h,\mathfrak h')$ in the available
block spans such that
\begin{equation}\label{eq:lm:type-correction}
 L(\mathfrak h,\mathfrak h')+
 L(\mathfrak h',\mathfrak h)^t=b(\mathfrak h,\mathfrak h').
\end{equation}
For distinct types, choose either direction using
Equation~\eqref{eq:lm:symmetry-defect}. For equal types, the right
side is alternating. It is a matrix plus its transpose in the full
span; in the rotation span use $R+R^{-1}=1$ after multiplication by
$J$. On a simple line it is zero. Choose $L=0$ when one type is
primary and trivial at $S$, since then its support pairings vanish.
Define, on every pair including the diagonal,
\[
 M_{qp}=B_{qp}-L(\mathfrak h_q,\mathfrak h_p).
\]
The matrix $M$ is symmetric. Lemma~\ref{lem:lm:arithmetic-realization}
allows zero cross-$M$ entries between prescribed blocks independently
of their internal entries.

A \emph{stationary environment} consists of the fixed coordinates
outside the variable blocks; its dimension is the number of those
coordinates.

\begin{proposition}[Nullity forces Selmer dimension]
\label{prop:selmer-matrix-bound}
Fix the curve, support types, a bound for the number of additional
constrained places, and a bound for the dimension of a stationary
environment. Partition a finite configuration into blocks with zero
cross-$M$ entries, and denote their matrices by $M_i$. There is a
constant depending only on these fixed data such that
\begin{equation}\label{eq:lm:selmer-nullity}
 \dim_{\F_2}\Sel_2(E^{(h)}/\Q)
 \geq\sum_i\dim_{\F_2}\ker M_i-C.
\end{equation}
At a unit coefficient symbol the sum of nullities, and hence the
number of singular blocks, is therefore uniformly bounded. The bound
works for both coefficient sources and is independent of the moving
prime $\ell$.
\end{proposition}

\begin{proof}
On $\bigoplus_i\ker M_i$, impose
\begin{equation}\label{eq:lm:type-sums}
 \sum_{p:\,\mathfrak h_p=\mathfrak h}x_p=0
 \quad\text{for each type }\mathfrak h,
\end{equation}
using the fixed orientation and basis to identify coordinates of the
same type. These equations have codimension at most
$\sum_{\mathfrak h}\dim V_{\mathfrak h}$. They kill the correction
in every row, since that correction is
\[
 \sum_{\mathfrak h}L(\mathfrak h_q,\mathfrak h)
       \sum_{p:\,\mathfrak h_p=\mathfrak h}x_p.
\]
Thus $Bx=0$ at all block positions. Set the environment coordinates
to zero and impose its row conditions, losing at most its dimension.
The Selmer conditions at $S$ and the additional constrained places
have bounded total codimension in local $H^1$. At every remaining
good place the classes satisfy the unramified Kummer condition; at a
root prime local cohomology is zero. Their sums are therefore Selmer
classes, and independent inertia coordinates make the map injective.
This proves Equation~\eqref{eq:lm:selmer-nullity}.

At a unit symbol, Proposition~\ref{prop:unit-symbol-detection} bounds
the Selmer dimension, giving the remaining assertions. The type
constraints use only the fixed $S$; a condition at the single moving
prime contributes a uniformly bounded local dimension.
\end{proof}

\subsection{Three states on two coordinates}

We have obtained two kinds of information from an arithmetic prime
configuration: its coefficient symbol and a matrix whose nullity
bounds its Selmer dimension from below. To compare configurations,
we now vary exactly two coordinates at a time. Both the symbol and
the determinant obey one common three-state relation. They remain
different functions; the relation gives linearity, not their equality.
The bounded number of singular blocks at a unit symbol will be the
additional input used in the next section.

A \emph{site} consists of either one split vertex or an ordered pair
of simple vertices, and always supplies two matrix coordinates.
At a split site use the three states
\[
 \text{absent},\qquad
 \text{present with diagonal }Q,\qquad
 \text{present with diagonal }Q+J.
\]
The rational Legendre bit of the aggregate root selects the two
present diagonals. At a simple-pair site use absence of both vertices
and two present variants differing in their nonmatching cross bit.
Replace the second prime by a prime with the same template, self data,
norm data, and outside cross bits. Its matching bit with the first
prime must also change to keep the rational norm bit fixed; that bit
has zero matrix contribution. Thus the two present matrices again
differ by the symmetric toggle on their coordinate pair. Require all
alternatives to agree against alternatives at other sites. These are
finite simultaneous prescriptions allowed by
Lemma~\ref{lem:lm:arithmetic-realization}.

In comparing states, keep fixed the total parameter label at the fixed
support: the product of the labels of the active sites and the aggregate
roots. The aggregate-root label itself changes when a site is deleted,
absorbing the deleted labels.

\begin{lemma}[The site relations]\label{lem:lm:site-relations}
Fix the total parameter label at the fixed support and the diagonal
choices at the other present sites. In each state choose aggregate
roots restoring this label, compatibly with an order-three residual
Frobenius, and preserving those diagonals. For either coefficient
symbol, its three values at one site sum to zero. The same holds for
$\det M$, with empty determinant one.
Consequently both functions extend linearly to the tensor product of
one copy of $\F_2^2$ per site, whose three nonzero vectors represent
the three states. These factors encode evaluations, not vectors of $V$.
\end{lemma}

\begin{proof}
For a split site this is the root-deletion rule of
Theorem~\ref{thm:symbol-rules}. Deletion transfers the removed
vertex's rational labels to the aggregate root. Signed prime
discriminants preserve the unit twist, and hence the diagonal,
at every remaining prime.

For a simple-pair site, omit the varied second prime from its
opposite-parity trace test. In an extension of $G$ by an elementary
kernel, two lifts of its transposition $c$ are conjugate by the kernel
exactly when their difference lies in the image of $1+c$. Dually,
this means agreement on the $c$-fixed radical squareclasses.
Besides fixed-support radicals, these are generated by products of
the orbit generators $g\pi_p$ over $c$-orbits: their off-support
divisors are independent and conjugation permutes them. Singleton
orbits test matching bits. A two-element orbit has product in $F^c$;
its symbol at the varying degree-one prime is, by the residue norm,
the symbol of one factor in $F$, namely the nonmatching bit.

The variants agree on support radicals and the orbit tests for every
other vertex. At the first prime they change both matching and
nonmatching bits. Modulo the image of $1+c$, their difference is
therefore the product of inertia in the matching transposition
orientation and inertia in the other orientation. Apply these
substitutions successively. The matching-orientation derivative is zero;
the different-transposition derivative deletes the pair by
Theorem~\ref{thm:symbol-rules}. Multiplication
by inertia leaves the residual transposition unchanged. The intermediate
Frobenius test is realized by Chebotarev; it need not be a site variant.
The resulting root has rotation image and transfers the pair's rational
labels to the aggregate. Labels elsewhere are unchanged. At simple
vertices root probes are immaterial; at split vertices the common
rational labels preserve the diagonals. Matched actual probes, or their
aggregate, are available because their discriminant character is trivial.
This proves the symbol relation.

For the determinant, reverse every permutation cycle of length at
least three. The resulting terms cancel in characteristic two, leaving
only involutions. Toggling the symmetric entries $(i,j),(j,i)$ changes
exactly the terms containing $(ij)$, whose sum is the determinant with
coordinates $i,j$ deleted. This includes the split $J$-shift. Finally,
the three-state relation in each factor is precisely linearity on its
three nonzero binary vectors.
\end{proof}

\begin{lemma}[Binary tensor span]\label{lem:binary-tensor-span}
Let $W=\bigotimes_{i=1}^rW_i$, where each $W_i$ has dimension two
over $\F_2$, and let $0\ne f\in W^*$. The nonzero pure tensors
annihilated by $f$ span $\ker f$. Thus if pure tensors $e,t$ satisfy
$f(e)=f(t)=1$, then $t$ is $e$ plus a sum of pure zeros of $f$.
\end{lemma}

\begin{proof}
It suffices to show that a functional $g$ vanishing on the pure zeros
of $f$ is either $0$ or $f$. Induct on $r$, the case $r=0$ being
immediate. Slice at the three nonzero vectors of the last factor,
numbering so that $f_3=f_1+f_2$ and $g_3=g_1+g_2$.
If $f_i=0$, then $g_i=0$ because the smaller pure tensors span.
Otherwise induction gives $g_i=\epsilon_i f_i$ for a binary scalar
$\epsilon_i$. If all three $f_i$ are nonzero, $f_1,f_2$ are independent,
so the sum identities force all $\epsilon_i$ to agree. If one $f_i$
vanishes, the other two, and their $g_i$, agree. All three cannot vanish.
Hence $g=0$ or $g=f$.
\end{proof}

Take $f=\det M$ and let $e$ be the all-absent state. Every invertible
state is the empty state plus a sum of singular pure states.
Moreover, on any selected set of active sites, some choice of pair
toggles makes $M$ invertible: the iterated determinant difference over
all selected pairs is the empty determinant, one. These facts turn the
Selmer nullity bound into a constraint on symbol evaluations.

\subsection{Making one address singular}

The final construction singles out one nonzero vector of a binary
cube. Every position is activated by a linear form, as required for
interpolation; auxiliary split sites repair the matrix at all other
nonzero vectors.

\begin{lemma}[A block address program]\label{lem:lm:address-program}
Let $M_*$ be the singular matrix of a finite site state, and let $D$
be a sum of its pair toggles such that $N_*=M_*+D$ is invertible.
For $b\geq2$ and $0\ne x_*\in\F_2^b$, one can add finitely many
primary split sites, each activated by a linear form, so that:
\begin{itemize}
\item all sites present in the original state have activation $l_0$,
with $l_0(x_*)=1$;
\item no site is active at zero, and at $x_*$ precisely the original
state $M_*$ remains;
\item the matrix is invertible at every other nonzero address;
\item the active auxiliary matrix is always invertible, and its Schur
complement contributes $D$ exactly when the original sites are active
and $x\ne x_*$.
\end{itemize}
All entries use the freedoms of
Lemma~\ref{lem:lm:arithmetic-realization}.
\end{lemma}

\begin{proof}
Choose a basis $y=(y_1,\ldots,y_{b-1})$ for the forms annihilating
$x_*$ and complete it by $l_0$. Then $y=0$ exactly at $0,x_*$.
For each required toggle and each nonzero $\eta\in\F_2^{b-1}$,
choose a basis of forms all equal to one at $\eta$. Such a basis is
obtained from a form $u$ with $u(\eta)=1$ by adding $u$ to a basis
of the annihilator of $\eta$. Its unique simultaneous all-one
solution is $\eta$.

Construct a path with $r=b-1$ steps, using two new split positions at
each step, both activated by the corresponding basis form. In scalar
block notation, where a scalar denotes that scalar times $J$, put
\begin{equation}\label{eq:lm:path-matrix}
 P=\begin{pmatrix}0&T\\T^t&B_e\end{pmatrix},
 \qquad T=I+\text{lower consecutive shift}.
\end{equation}
The near and far endpoints are steps $1$ and $r$. The matrix $B_e$
has one nonzero diagonal entry at the far endpoint. Restrict both sides
to active steps, retaining only the original consecutive edges. The
restricted $T$ is triangular
with diagonal one, so $P$ is invertible, with scalar inverse
\[
 P^{-1}=\begin{pmatrix}
 T^{-t}B_eT^{-1}&T^{-t}\\T^{-1}&0
 \end{pmatrix}.
\]
Couple the addressed coordinate pair by $J$ to the first side at the
near endpoint. Such a pairing is available also for two simple lines.
Solving $Tv=e_{\rm near}$ propagates one to the far endpoint exactly
when every step is active; a missing step interrupts the consecutive
shift. The Schur contribution is consequently $J$ on a complete path
and zero otherwise, including when the near endpoint is absent.

For each nonzero $y$, exactly one path, the one indexed by $\eta=y$,
is complete. Their direct sum supplies the toggle when $y\ne0$ and
supplies zero when $y=0$. Repeat for each toggle in $D$, with zero
mutual cross-$M$ entries. The actual matrices are the scalar matrices
tensored with $J$, so the same inverse calculation gives the asserted
binary Schur complements. All blocks lie in the permitted spans, and
primary diagonals are $0$ or $J$.

At $x_*$ the auxiliary sites are absent and $l_0=1$, leaving $M_*$.
At any other nonzero address with $l_0=1$, eliminating the auxiliary
sites leaves $M_*+D=N_*$. If $l_0=0$, only the invertible auxiliary
matrix remains. Every activation is linear and vanishes at zero.
\end{proof}
